\shead{Collecting and presenting data}
% ---- chapter 11 helpers: bar chart, pie chart with legend, pictograph icons ----
\newcommand{\stpBars}[6]{% #1 cm per unit, #2 ymax, #3 ystep, #4 list name/height/shown, #5 y caption, #6 x caption
\begin{tikzpicture}[font=\footnotesize,y=#1cm,x=1cm,>=stealth]
\foreach \n/\v/\d [count=\i] in {#4}{\xdef\stpN{\i}}
\pgfmathsetmacro\stpW{\stpN*1.2+0.1}
\foreach \t in {0,#3,...,#2}{\draw[Line] (0,\t)--(\stpW,\t); \node[left,text=Slate] at (0,\t){$\t$};}
\foreach \n/\v/\d [count=\i] in {#4}{
  \pgfmathsetmacro\stpC{\i*1.2-0.5}
  \fill[tmBlue!65] (\stpC-0.3,0) rectangle (\stpC+0.3,\v);
  \node[above,font=\footnotesize\bfseries,text=Ink] at (\stpC,\v){$\d$};
  \node[below,text=Ink] at (\stpC,0){\n};}
\draw[->,line width=0.8pt] (0,0)--([yshift=0.55cm]0,#2);
\draw[->,line width=0.8pt] (0,0)--(\stpW+0.35,0);
\node[above,text=Slate] at ([yshift=0.55cm]0,#2){#5};
\node[right,text=Slate] at (\stpW+0.35,0){#6};
\end{tikzpicture}}
\newcommand{\stpPie}[2][1.5]{% #2 list value/text/colour/legend name (value in percent)
\begin{tikzpicture}[font=\footnotesize]
\def\stpA{90}
\foreach \v/\t/\c/\n [count=\i] in {#2}{
  \pgfmathsetmacro\stpB{\stpA-3.6*\v}
  \pgfmathsetmacro\stpM{(\stpA+\stpB)/2}
  \pgfmathsetmacro\stpR{0.64*#1}
  \fill[\c] (0,0) -- (\stpA:#1) arc[start angle=\stpA,end angle=\stpB,radius=#1] -- cycle;
  \draw[white,line width=1.2pt] (0,0) -- (\stpA:#1);
  \node[text=Ink,font=\footnotesize\bfseries] at (\stpM:\stpR) {\t};
  \fill[\c] (#1+0.7,#1-0.55*\i+0.1) rectangle ++(0.32,0.32);
  \node[anchor=west,text=Ink,font=\small] at (#1+1.1,#1-0.55*\i+0.26) {\n};
  \global\let\stpA\stpB}
\draw[Slate,line width=0.8pt] (0,0) circle (#1);
\end{tikzpicture}}
\newcommand{\stpIcons}[1]{\foreach \k in {1,...,#1}{\textcolor{tmBlue}{\faBook}\ }}
\begin{knowbox}
To compare, judge and decide, people \textbf{collect information} (data) and then \textbf{organise} it. Sometimes we ask everybody; sometimes we ask only a typical group.\par\smallskip
Data are usually shown as:
\begin{itemize}[leftmargin=1.4em,itemsep=3pt]
\item a \textbf{table} -- each value is written in its own cell;
\item a \textbf{pictograph} -- each picture stands for the same number of things;
\item a \textbf{bar chart} -- the height of each bar, read on the vertical scale, gives the value;
\item a \textbf{pie chart} -- the whole circle is the whole total ($100\%$); each slice is a percentage of the total. Half the circle is $50\%$, a quarter is $25\%$.
\end{itemize}
From a pie chart: \quad \textbf{part $=$ total $:100\times$ percentage}; \quad \textbf{total $=$ part $:$ percentage $\times100$}.
\end{knowbox}
\shead{The mean (average)}
\begin{knowbox}
\textbf{Mean} $=$ (sum of the values) $:$ (number of values). \qquad \textbf{Sum} $=$ mean $\times$ number of values.
\begin{itemize}[leftmargin=1.4em,itemsep=3pt]
\item The mean always lies between the smallest and the largest value.
\item To find a \emph{missing value}: work out the sum from the mean, then subtract the values you know.
\item If one value grows by some amount, the sum grows by that amount and the mean grows by (that amount) $:$ (number of values).
\end{itemize}
\end{knowbox}
\shead{Chance and counting outcomes}
\begin{knowbox}
When we repeat an experiment (tossing a coin, rolling a dice, spinning a spinner, picking a ball without looking), we record how many times each outcome happens.
\begin{center}
\textbf{ratio describing how often an outcome happened} $=$ $\dfrac{\text{number of times it happened}}{\text{total number of trials}}$
\end{center}
For example, a dice is rolled $20$ times and the face with $6$ dots comes up $3$ times: the ratio is $\dfrac{3}{20}$. If we roll more times, this ratio may change.\par\smallskip
\textbf{Counting outcomes.} List every possible result in an orderly way (a list, a table or a tree). A coin has $2$ outcomes, a dice has $6$; a coin \emph{and} a dice together have $2\times6=12$.\par
When all outcomes are \emph{equally likely}, the \textbf{chance} of an event $=$ (number of outcomes that give the event) $:$ (number of all outcomes). An event that always happens is \emph{certain}; one that can never happen is \emph{impossible}.
\end{knowbox}
\gap
\begin{exbox}{\ExWord\ 1}
Lan writes down how many pages she reads each day of one week:
\begin{center}
\begin{tabular}{|l|c|c|c|c|c|c|c|}
\hline
\textbf{Day} & Mon & Tue & Wed & Thu & Fri & Sat & Sun\\\hline
\textbf{Pages} & $18$ & $25$ & $20$ & $32$ & $15$ & $30$ & $28$\\\hline
\end{tabular}
\end{center}
\begin{cols2}
\pa{a} On which day did she read most? Least? & \pa{b} How many pages did she read per day on average?\\
\multicolumn{2}{@{}l@{}}{\pa{c} On how many days did she read more than her average?}
\end{cols2}
\guide
\pa{a} Most on Thursday ($32$ pages); least on Friday ($15$ pages).\par
\pa{b} Pages in the week: $18+25+20+32+15+30+28=168$.\par Mean per day: $168:7=24$ pages.\par
\pa{c} More than $24$: Tuesday, Thursday, Saturday and Sunday, i.e.\ $4$ days.
\dapso{\pa{a} Thursday; Friday. \pa{b} $24$ pages. \pa{c} $4$ days.}
\end{exbox}
\gap
\begin{exbox}{\ExWord\ 2}
The bar chart shows how many books five classes gave to the school library.
\begin{center}
\stpBars{0.05}{80}{10}{5A/45/45,5B/60/60,5C/35/35,5D/70/70,5E/40/40}{(books)}{(class)}
\end{center}
\begin{cols2}
\pa{a} How many books did the five classes give? & \pa{b} How many books did a class give on average?\\
\multicolumn{2}{@{}l@{}}{\pa{c} How many more books did the class that gave most give than the class that gave least?}
\end{cols2}
\guide
\pa{a} $45+60+35+70+40=250$ books.\par
\pa{b} $250:5=50$ books.\par
\pa{c} Most: 5D ($70$); least: 5C ($35$).\par Difference: $70-35=35$ books.
\dapso{\pa{a} $250$ books. \pa{b} $50$ books. \pa{c} $35$ books.}
\end{exbox}
\gap
\begin{exbox}{\ExWord\ 3}
$400$ pupils were asked about their favourite fruit. The results are shown in the pie chart.
\begin{center}
\stpPie{35/35\%/tmBlue!55/Mango,20/20\%/tmOrange!60/Banana,25/25\%/tmGreen!50/Watermelon,20/?/tmPurple!40/Orange}
\end{center}
\pa{a} What percentage of the pupils chose oranges?\par
\pa{b} How many pupils chose mangoes?\par
\pa{c} How many more pupils chose mangoes than bananas?
\guide
\pa{a} The whole circle is $100\%$: $100\%-35\%-20\%-25\%=20\%$.\par
\pa{b} $400:100\times35=140$ pupils.\par
\pa{c} \emph{Method 1.} Bananas: $400:100\times20=80$ pupils; \ $140-80=60$ pupils.\par
\emph{Method 2.} Mangoes are $35\%-20\%=15\%$ more; \ $400:100\times15=60$ pupils.
\dapso{\pa{a} $20\%$. \pa{b} $140$ pupils. \pa{c} $60$ pupils.}
\end{exbox}
\gap
\begin{exbox}{\ExWord\ 4}
Minh rolls a dice $40$ times and records the results.
\begin{center}
\begin{tabular}{|l|c|c|c|c|c|c|}
\hline
\textbf{Outcome} & $1$ dot & $2$ dots & $3$ dots & $4$ dots & $5$ dots & $6$ dots\\\hline
\textbf{Times} & $6$ & $8$ & $5$ & $7$ & $9$ & $5$\\\hline
\end{tabular}
\end{center}
Write the ratio describing how often each event happened, compared with the number of rolls:\par
\pa{a} $6$ dots; \quad \pa{b} an even number of dots; \quad \pa{c} at least $5$ dots.
\guide
\pa{a} $6$ dots came up $5$ times: $\dfrac{5}{40}=\dfrac{1}{8}$.\par
\pa{b} Even ($2$, $4$, $6$ dots): $8+7+5=20$ times; \ $\dfrac{20}{40}=\dfrac12$.\par
\pa{c} At least $5$ dots ($5$ or $6$ dots): $9+5=14$ times; \ $\dfrac{14}{40}=\dfrac{7}{20}$.
\note{This does not mean the same thing will happen next time: with more rolls the ratios change.}
\dapso{\pa{a} $\dfrac18$. \pa{b} $\dfrac12$. \pa{c} $\dfrac{7}{20}$.}
\end{exbox}
\gap
\begin{exbox}{\ExWord\ 5}
Box I holds three cards numbered $1$, $2$, $3$; box II holds two cards numbered $4$, $5$. Without looking, Hoa takes one card from each box and adds the two numbers.\par
\pa{a} List all possible outcomes. \quad \pa{b} What is the chance that the sum is $7$?
\guide
\pa{a} Make a table (row: box I, column: box II):
\begin{center}
\begin{tabular}{|c|c|c|}
\hline
$+$ & $4$ & $5$\\\hline
$1$ & $5$ & $6$\\\hline
$2$ & $6$ & $7$\\\hline
$3$ & $7$ & $8$\\\hline
\end{tabular}
\end{center}
There are $3\times2=6$ equally likely outcomes: $1+4,\ 1+5,\ 2+4,\ 2+5,\ 3+4,\ 3+5$.\par
\pa{b} The sum is $7$ in $2$ outcomes ($2+5$ and $3+4$).\par Chance: $\dfrac26=\dfrac13$.
\dapso{\pa{a} $6$ outcomes. \pa{b} $\dfrac13$.}
\end{exbox}
\dividerbar{Practice}
\begin{probbox}{1}{\relax}
The noon temperature in a town on five days was:
\begin{center}
\begin{tabular}{|l|c|c|c|c|c|}
\hline
\textbf{Day} & Mon & Tue & Wed & Thu & Fri\\\hline
\textbf{Temperature ($^\circ$C)} & $31$ & $29$ & $33$ & $30$ & $32$\\\hline
\end{tabular}
\end{center}
\pa{a} Which day was hottest, which coolest? \quad \pa{b} Find the mean noon temperature.
\end{probbox}
\probsol{
\pa{a} Hottest: Wednesday ($33^\circ$C); coolest: Tuesday ($29^\circ$C).\par
\pa{b} $(31+29+33+30+32):5=155:5=31^\circ$C.
\dapso{\pa{a} Wednesday; Tuesday. \pa{b} $31^\circ$C.}}
\gap
\begin{probbox}{2}{\relax}
Nam's marks in six maths tests were $7,\ 9,\ 8,\ 10,\ 6,\ 8$. What is his mean mark?
\end{probbox}
\probsol{
Sum: $7+9+8+10+6+8=48$.\par Mean: $48:6=8$.
\dapso{$8$.}}
\gap
\begin{probbox}{3}{\relax}
Four friends are $135$ cm, $142$ cm, $138$ cm and $145$ cm tall. Find their mean height.
\end{probbox}
\probsol{
$(135+142+138+145):4=560:4=140$ cm.
\dapso{$140$ cm.}}
\gap
\begin{probbox}{4}{\relax}
The pictograph shows how many books four groups read in May.
\begin{center}
\begin{tabular}{|l|l|}
\hline
Group 1 & \stpIcons{5}\\\hline
Group 2 & \stpIcons{3}\\\hline
Group 3 & \stpIcons{6}\\\hline
Group 4 & \stpIcons{4}\\\hline
\end{tabular}\par\smallskip
Each \textcolor{tmBlue}{\faBook} stands for $4$ books.
\end{center}
\pa{a} How many books did the four groups read? \quad \pa{b} How many books did a group read on average?
\end{probbox}
\probsol{
Number of pictures: $5+3+6+4=18$.\par
\pa{a} Books: $4\times18=72$.\par
\pa{b} Mean: $72:4=18$ books.
\dapso{\pa{a} $72$ books. \pa{b} $18$ books.}}
\gap
\begin{probbox}{5}{\relax}
A shop sold $24$ bicycles in January, $30$ in February, $27$ in March and $39$ in April.\par
\pa{a} How many bicycles were sold in the four months? \quad \pa{b} How many a month on average?\par
\pa{c} In how many months did the shop sell more than its average?
\end{probbox}
\probsol{
\pa{a} $24+30+27+39=120$ bicycles.\par
\pa{b} $120:4=30$ bicycles.\par
\pa{c} Only April ($39>30$; February's $30$ is equal to the mean, not more): $1$ month.
\dapso{\pa{a} $120$. \pa{b} $30$. \pa{c} $1$ month.}}
\gap
\begin{probbox}{6}{\relax}
The chart shows the number of pupils in each grade of a primary school.
\begin{center}
\stpBars{0.025}{150}{25}{Grade 1/120/120,Grade 2/135/135,Grade 3/110/110,Grade 4/140/140,Grade 5/125/125}{(pupils)}{}
\end{center}
\pa{a} Which grade has the most pupils? The fewest?\par
\pa{b} How many pupils does the school have?\par
\pa{c} How many more pupils does Grade 4 have than Grade 3?
\end{probbox}
\probsol{
\pa{a} Most: Grade 4 ($140$); fewest: Grade 3 ($110$).\par
\pa{b} $120+135+110+140+125=630$ pupils.\par
\pa{c} $140-110=30$ pupils.
\dapso{\pa{a} Grade 4; Grade 3. \pa{b} $630$ pupils. \pa{c} $30$ pupils.}}
\gap
\begin{probbox}{7}{\relax}
The chart shows the rainfall in a town from May to August.
\begin{center}
\stpBars{0.014}{250}{50}{May/80/80,June/150/150,July/210/210,August/120/120}{(mm)}{(month)}
\end{center}
Find the total rainfall and the mean rainfall per month.
\end{probbox}
\probsol{
Total: $80+150+210+120=560$ mm.\par Mean: $560:4=140$ mm.
\dapso{$560$ mm; $140$ mm.}}
\gap
\begin{probbox}{8}{\relax}
Class 5B has $40$ pupils. The pie chart shows how they come to school.
\begin{center}
\stpPie{30/30\%/tmBlue!55/On foot,45/45\%/tmOrange!60/By bicycle,25/25\%/tmGreen!50/By motorbike with parents}
\end{center}
\pa{a} How many pupils walk to school?\par
\pa{b} How many more pupils come by bicycle than by motorbike?
\end{probbox}
\probsol{
\pa{a} $40:100\times30=12$ pupils.\par
\pa{b} By bicycle: $40:100\times45=18$; by motorbike: $40:100\times25=10$.\par Difference: $18-10=8$ pupils.
\dapso{\pa{a} $12$ pupils. \pa{b} $8$ pupils.}}
\gap
\begin{probbox}{9}{\relax}
$200$ pupils chose their favourite sport.
\begin{center}
\stpPie{40/40\%/tmBlue!55/Football,25/25\%/tmOrange!60/Badminton,20/20\%/tmGreen!50/Swimming,15/15\%/tmPurple!40/Chess}
\end{center}
\pa{a} How many pupils chose football? \quad \pa{b} chess?\par
\pa{c} How many pupils chose football or swimming?
\end{probbox}
\probsol{
\pa{a} $200:100\times40=80$ pupils.\par
\pa{b} $200:100\times15=30$ pupils.\par
\pa{c} $40\%+20\%=60\%$; \ $200:100\times60=120$ pupils.
\dapso{\pa{a} $80$. \pa{b} $30$. \pa{c} $120$ pupils.}}
\gap
\begin{probbox}{10}{\relax}
$300$ pupils each joined one after-school club.
\begin{center}
\stpPie{30/30\%/tmBlue!55/Art,25/25\%/tmOrange!60/Music,25/?/tmGreen!50/Science,20/20\%/tmPurple!40/Sport}
\end{center}
What percentage of the pupils joined the Science club, and how many pupils is that?
\end{probbox}
\probsol{
Science: $100\%-30\%-25\%-20\%=25\%$.\par Pupils: $300:100\times25=75$.
\dapso{$25\%$; $75$ pupils.}}
\gap
\begin{probbox}{11}{\relax}
A family spends $12\,000\,000$ dong a month as the chart shows.
\begin{center}
\stpPie{50/50\%/tmBlue!55/Food,25/25\%/tmOrange!60/School,15/15\%/tmGreen!50/Savings,10/10\%/tmPurple!40/Other}
\end{center}
How much money goes on food? How much is saved?
\end{probbox}
\probsol{
Food is half: $12\,000\,000:2=6\,000\,000$ dong.\par
Savings: $12\,000\,000:100\times15=1\,800\,000$ dong.
\dapso{Food $6\,000\,000$ dong; savings $1\,800\,000$ dong.}}
\gap
\begin{probbox}{12}{\relax}
In a survey, $36$ pupils chose mango as their favourite fruit. They are $30\%$ of all the pupils asked. How many pupils were asked?
\end{probbox}
\probsol{
$30\%$ of the pupils are $36$ pupils, so all the pupils ($100\%$): $36:30\times100=120$.
\dapso{$120$ pupils.}}
\gap
\begin{probbox}{13}{\relax}
The mean of five numbers is $18$. Four of them are $15$, $20$, $17$ and $22$. Find the fifth number.
\end{probbox}
\probsol{
Sum of the five numbers: $18\times5=90$.\par Sum of the four: $15+20+17+22=74$.\par Fifth number: $90-74=16$.
\dapso{$16$.}}
\gap
\begin{probbox}{14}{\relax}
An's mean mark in three tests is $8$. After a fourth test his mean mark is $8.5$. What did he get in the fourth test?
\end{probbox}
\probsol{
Sum of three marks: $8\times3=24$.\par Sum of four marks: $8.5\times4=34$.\par Fourth mark: $34-24=10$.
\dapso{$10$.}}
\gap
\begin{probbox}{15}{\relax}
Three bags of rice weigh $25$ kg on average. A fourth bag of $33$ kg is added. What is the mean weight of the four bags?
\end{probbox}
\probsol{
Three bags: $25\times3=75$ kg.\par Four bags: $75+33=108$ kg.\par Mean: $108:4=27$ kg.
\dapso{$27$ kg.}}
\gap
\begin{probbox}{16}{\relax}
Tú tosses a coin $50$ times; it lands heads up $28$ times. Write the ratio describing how often it landed heads, and how often tails, compared with the number of tosses.
\end{probbox}
\probsol{
Heads: $\dfrac{28}{50}=\dfrac{14}{25}$.\par Tails: $50-28=22$ times, $\dfrac{22}{50}=\dfrac{11}{25}$.
\dapso{Heads $\dfrac{14}{25}$; tails $\dfrac{11}{25}$.}}
\gap
\begin{probbox}{17}{\relax}
A spinner with red, blue and yellow parts is spun $60$ times: red $21$ times, blue $15$ times, yellow $24$ times.\par
\pa{a} Write the ratio for each colour. \quad \pa{b} Write the ratio for ``not red''.
\end{probbox}
\probsol{
\pa{a} Red: $\dfrac{21}{60}=\dfrac{7}{20}$; \ blue: $\dfrac{15}{60}=\dfrac14$; \ yellow: $\dfrac{24}{60}=\dfrac25$.\par
\pa{b} Not red: $15+24=39$ times; \ $\dfrac{39}{60}=\dfrac{13}{20}$.
\dapso{\pa{a} $\dfrac{7}{20}$; $\dfrac14$; $\dfrac25$. \pa{b} $\dfrac{13}{20}$.}}
\gap
\begin{probbox}{18}{\relax}
A dice is rolled $30$ times:
\begin{center}
\begin{tabular}{|l|c|c|c|c|c|c|}
\hline
\textbf{Outcome} & $1$ dot & $2$ dots & $3$ dots & $4$ dots & $5$ dots & $6$ dots\\\hline
\textbf{Times} & $4$ & $6$ & $5$ & $3$ & $7$ & $5$\\\hline
\end{tabular}
\end{center}
Write the ratio for: \pa{a} an odd number of dots; \quad \pa{b} more than $4$ dots.
\end{probbox}
\probsol{
\pa{a} Odd ($1$, $3$, $5$): $4+5+7=16$ times; \ $\dfrac{16}{30}=\dfrac{8}{15}$.\par
\pa{b} More than $4$ ($5$, $6$): $7+5=12$ times; \ $\dfrac{12}{30}=\dfrac25$.
\dapso{\pa{a} $\dfrac{8}{15}$. \pa{b} $\dfrac25$.}}
\gap
\begin{probbox}{19}{\relax}
Bình rolls a dice many times. The face with $6$ dots comes up $9$ times, and the ratio for $6$ dots is $\dfrac{3}{20}$. How many times did Bình roll the dice? How many rolls did not show $6$ dots?
\end{probbox}
\probsol{
$9$ times is $3$ parts out of $20$, so one part is $9:3=3$ rolls.\par Rolls: $3\times20=60$.\par Not $6$ dots: $60-9=51$.
\dapso{$60$ rolls; $51$ rolls.}}
\gap
\begin{probbox}{20}{\relax}
A bag holds red and blue balls. Mai takes a ball without looking, writes its colour and puts it back. She does this $25$ times, and the ratio for red is $\dfrac25$. How many times did she get red, and how many times blue?
\end{probbox}
\probsol{
Red: $25\times\dfrac25=10$ times.\par Blue: $25-10=15$ times.
\dapso{Red $10$ times, blue $15$ times.}}
\gap
\begin{probbox}{21}{\relax}
Two coins are tossed at the same time.\par
\pa{a} List all possible outcomes. \quad \pa{b} In how many outcomes is there one head and one tail?
\end{probbox}
\probsol{
Write H for heads and T for tails (first coin, then second coin).\par
\pa{a} HH, HT, TH, TT: $4$ outcomes.\par
\pa{b} HT and TH: $2$ outcomes.
\dapso{\pa{a} $4$ outcomes. \pa{b} $2$ outcomes.}}
\gap
\begin{probbox}{22}{\relax}
A school canteen offers $3$ main dishes (rice, noodles, bread) and $2$ drinks (milk, juice). A lunch is one main dish and one drink. How many different lunches are there? List them.
\end{probbox}
\probsol{
Each dish goes with $2$ drinks: $3\times2=6$ lunches.\par
Rice--milk, rice--juice, noodles--milk, noodles--juice, bread--milk, bread--juice.
\dapso{$6$ lunches.}}
\gap
\begin{probbox}{23}{\relax}
Using the digits $1$, $4$, $7$, how many two-digit numbers with two different digits can be written? How many of them are even?
\end{probbox}
\probsol{
Numbers: $14,\ 17,\ 41,\ 47,\ 71,\ 74$: $6$ numbers ($3$ choices for the tens digit, then $2$ for the ones).\par
Even numbers end in $4$: $14$ and $74$, i.e.\ $2$ numbers.
\dapso{$6$ numbers; $2$ even.}}
\gap
\begin{probbox}{24}{\relax}
A bag holds $3$ red, $5$ blue and $2$ yellow balls of the same size. One ball is taken without looking.\par
\pa{a} Which colour is most likely? \quad \pa{b} What is the chance of a blue ball?\par
\pa{c} Is a green ball certain, possible or impossible?
\end{probbox}
\probsol{
\pa{a} Blue: there are most blue balls ($5$).\par
\pa{b} $10$ balls, $5$ blue: $\dfrac{5}{10}=\dfrac12$.\par
\pa{c} There are no green balls: impossible.
\dapso{\pa{a} Blue. \pa{b} $\dfrac12$. \pa{c} Impossible.}}
\gap
\begin{probbox}{25}{\relax}
A football team scored $2,\ 0,\ 3,\ 1,\ 4,\ 2,\ 1,\ 3$ goals in eight matches. Find the mean number of goals per match. In how many matches did the team score more than the mean?
\end{probbox}
\probsol{
Goals: $2+0+3+1+4+2+1+3=16$.\par Mean: $16:8=2$ goals.\par More than $2$: the matches with $3,\ 4,\ 3$ goals, i.e.\ $3$ matches.
\dapso{$2$ goals; $3$ matches.}}
\gap
\begin{probbox}{26}{\relax}
The chart shows the rice harvest of a farm (in tonnes).
\begin{center}
\stpBars{0.06}{60}{10}{2021/42/42,2022/45/45,2023/48/48,2024/50/50,2025/55/55}{(tonnes)}{(year)}
\end{center}
\pa{a} Find the mean harvest per year. \quad \pa{b} How much more was harvested in 2025 than in 2021?
\end{probbox}
\probsol{
\pa{a} $(42+45+48+50+55):5=240:5=48$ tonnes.\par
\pa{b} $55-42=13$ tonnes.
\dapso{\pa{a} $48$ tonnes. \pa{b} $13$ tonnes.}}
\gap
\begin{probbox}{27}{\relax}
A library has $1\,200$ books.
\begin{center}
\stpPie{35/35\%/tmBlue!55/Comics,25/25\%/tmOrange!60/Science,30/30\%/tmGreen!50/Stories,10/10\%/tmPurple!40/Other}
\end{center}
\pa{a} How many science books are there? \quad \pa{b} How many more story books than ``other'' books?
\end{probbox}
\probsol{
\pa{a} $1\,200:100\times25=300$ books.\par
\pa{b} Stories are $30\%-10\%=20\%$ more: $1\,200:100\times20=240$ books.
\dapso{\pa{a} $300$ books. \pa{b} $240$ books.}}
\gap
\begin{probbox}{28}{\relax}
Hoa spent her $160\,000$ dong of pocket money as shown.
\begin{center}
\stpPie{50/50\%/tmBlue!55/Books,25/25\%/tmOrange!60/Savings,15/15\%/tmGreen!50/Snacks,10/10\%/tmPurple!40/Gifts}
\end{center}
\pa{a} How much did she save? \quad \pa{b} How much more did she spend on books than on snacks?
\end{probbox}
\probsol{
\pa{a} A quarter: $160\,000:4=40\,000$ dong.\par
\pa{b} Books: $80\,000$ dong; snacks: $160\,000:100\times15=24\,000$ dong.\par Difference: $80\,000-24\,000=56\,000$ dong.
\dapso{\pa{a} $40\,000$ dong. \pa{b} $56\,000$ dong.}}
\gap
\begin{probbox}{29}{\relax}
Five girls weigh $32$ kg on average and three boys weigh $36$ kg on average. Find the mean weight of all eight children.
\end{probbox}
\probsol{
Girls: $32\times5=160$ kg.\par Boys: $36\times3=108$ kg.\par Mean: $(160+108):8=268:8=33.5$ kg.
\dapso{$33.5$ kg.}}
\gap
\begin{probbox}{30}{\relax}
Lan tosses a coin several times. It lands tails up $14$ times, and the ratio for heads is $\dfrac35$. How many times did she toss the coin, and how many heads were there?
\end{probbox}
\probsol{
Ratio for tails: $1-\dfrac35=\dfrac25$.\par So $14$ times is $2$ parts out of $5$: one part is $7$ tosses.\par Tosses: $7\times5=35$; heads: $35-14=21$.
\dapso{$35$ tosses; $21$ heads.}}
\dividerbar{Challenge -- Extension}
\begin{probbox}{31}{\stars{1}}
The mean of six numbers is $15$. When one number is removed, the mean of the other five is $14$. Which number was removed?
\end{probbox}
\probsol{
Six numbers: $15\times6=90$.\par Five numbers: $14\times5=70$.\par Removed: $90-70=20$.
\dapso{$20$.}}
\gap
\begin{probbox}{32}{\stars{1}}
The mean age of $10$ players in a team is $30$. When the coach is counted too, the mean age is $31$. How old is the coach?
\end{probbox}
\probsol{
Players: $30\times10=300$ years.\par With the coach ($11$ people): $31\times11=341$ years.\par Coach: $341-300=41$.
\dapso{$41$ years old.}}
\gap
\begin{probbox}{33}{\stars{1}}
Hùng's mean mark in four tests is $7.5$. What must he get in the fifth test so that his mean mark becomes $8$?
\end{probbox}
\probsol{
Four tests: $7.5\times4=30$.\par Five tests: $8\times5=40$.\par Fifth test: $40-30=10$.
\dapso{$10$.}}
\gap
\begin{probbox}{34}{\stars{1}}
The mean of three numbers is $35$. The first is twice the second, and the second is twice the third. Find the three numbers.
\end{probbox}
\probsol{
Sum: $35\times3=105$.\par If the third number is $1$ part, the second is $2$ parts and the first is $4$ parts: $7$ parts.\par
One part: $105:7=15$.\par Numbers: $60$, $30$, $15$.
\dapso{$60$, $30$, $15$.}}
\gap
\begin{probbox}{35}{\stars{1}}
The mean of five consecutive even numbers is $36$. Find the largest of them.
\end{probbox}
\probsol{
For equally spaced numbers the mean is the middle number: $36$.\par The numbers are $32,\ 34,\ 36,\ 38,\ 40$.
\dapso{$40$.}}
\gap
\begin{probbox}{36}{\stars{1}}
The mean of four marks was worked out as $8.25$, but one mark was copied as $6$ instead of $9$. What is the correct mean?
\end{probbox}
\probsol{
Wrong sum: $8.25\times4=33$.\par Correct sum: $33-6+9=36$.\par Correct mean: $36:4=9$.
\dapso{$9$.}}
\gap
\begin{probbox}{37}{\stars{1}}
The chart shows tickets sold at a museum from Monday to Friday; the bar for Wednesday is missing. The mean is $160$ tickets a day. How many tickets were sold on Wednesday?
\begin{center}
\stpBars{0.016}{250}{50}{Mon/120/120,Tue/150/150,Wed/0/?,Thu/180/180,Fri/210/210}{(tickets)}{(day)}
\end{center}
\end{probbox}
\probsol{
Five days: $160\times5=800$ tickets.\par The other four days: $120+150+180+210=660$.\par Wednesday: $800-660=140$ tickets.
\dapso{$140$ tickets.}}
\gap
\begin{probbox}{38}{\stars{1}}
Each pupil of class 5A chose one sport. $12$ pupils chose basketball.
\begin{center}
\stpPie{40/40\%/tmBlue!55/Football,30/30\%/tmOrange!60/Basketball,20/20\%/tmGreen!50/Swimming,10/10\%/tmPurple!40/Chess}
\end{center}
\pa{a} How many pupils are there in class 5A? \quad \pa{b} How many chose football? Chess?
\end{probbox}
\probsol{
\pa{a} $12$ pupils are $30\%$: $12:30\times100=40$ pupils.\par
\pa{b} Football: $40:100\times40=16$; \ chess: $40:100\times10=4$.
\dapso{\pa{a} $40$ pupils. \pa{b} $16$; $4$.}}
\gap
\begin{probbox}{39}{\stars{2}}
$160$ pupils named their favourite pet: dog $40\%$, rabbit $15\%$, and the rest chose cat or fish. Twice as many pupils chose cats as chose fish. How many pupils chose cats?
\end{probbox}
\probsol{
Cat and fish: $100\%-40\%-15\%=45\%$.\par Fish $1$ part, cat $2$ parts: one part is $45\%:3=15\%$.\par Cat: $30\%$, i.e.\ $160:100\times30=48$ pupils.
\dapso{$48$ pupils.}}
\gap
\begin{probbox}{40}{\stars{1}}
In class 5A ($40$ pupils), $35\%$ like reading best. In class 5B ($36$ pupils), $50\%$ like reading best. Which class has more pupils who like reading best, and by how many?
\end{probbox}
\probsol{
5A: $40:100\times35=14$ pupils.\par 5B: $36:2=18$ pupils.\par 5B has $18-14=4$ more.
\dapso{Class 5B, by $4$ pupils.}}
\gap
\begin{probbox}{41}{\stars{1}}
Using the digits $0,\ 3,\ 5,\ 8$, how many two-digit numbers with different digits can be written? How many of them are divisible by $5$?
\end{probbox}
\probsol{
Tens digit: $3$ choices ($3$, $5$, $8$; not $0$); ones digit: $3$ choices (any other digit).\par Numbers: $3\times3=9$.\par
Divisible by $5$ means ending in $0$ or $5$: $30,\ 50,\ 80,\ 35,\ 85$: $5$ numbers.
\dapso{$9$ numbers; $5$ of them.}}
\gap
\begin{probbox}{42}{\stars{1}}
Six friends meet, and every two of them shake hands once. How many handshakes are there?
\end{probbox}
\probsol{
Each person shakes $5$ hands: $6\times5=30$, but each handshake is counted twice.\par Handshakes: $30:2=15$.
\dapso{$15$ handshakes.}}
\gap
\begin{probbox}{43}{\stars{2}}
Two dice are rolled together and the dots are added.\par
\pa{a} How many outcomes are there? \quad \pa{b} In how many outcomes is the sum $7$? What is the chance of a sum of $7$?\par
\pa{c} What is the chance of a sum of $12$?
\end{probbox}
\probsol{
\pa{a} Each of the $6$ faces of the first dice goes with $6$ faces of the second: $6\times6=36$ outcomes.\par
\pa{b} Sum $7$: $1+6,\ 2+5,\ 3+4,\ 4+3,\ 5+2,\ 6+1$: $6$ outcomes. Chance: $\dfrac{6}{36}=\dfrac16$.\par
\pa{c} Only $6+6$: chance $\dfrac{1}{36}$.
\dapso{\pa{a} $36$. \pa{b} $6$; $\dfrac16$. \pa{c} $\dfrac{1}{36}$.}}
\gap
\begin{probbox}{44}{\stars{1}}
A box holds $4$ red and $6$ blue balls. \pa{a} What is the chance of taking a red ball without looking?\par
\pa{b} How many red balls must be added so that the chance of red becomes $\dfrac12$?
\end{probbox}
\probsol{
\pa{a} $\dfrac{4}{10}=\dfrac25$.\par
\pa{b} The chance is $\dfrac12$ when there are as many red as blue balls: $6$ red.\par Add $6-4=2$ red balls.
\dapso{\pa{a} $\dfrac25$. \pa{b} $2$ red balls.}}
\gap
\begin{probbox}{45}{\stars{1}}
A bag holds $20$ red, blue and green balls. The chance of a red ball is $\dfrac14$ and the chance of a blue ball is $\dfrac25$. How many balls of each colour are there?
\end{probbox}
\probsol{
Red: $20\times\dfrac14=5$.\par Blue: $20\times\dfrac25=8$.\par Green: $20-5-8=7$.
\dapso{$5$ red, $8$ blue, $7$ green.}}
\gap
\begin{probbox}{46}{\stars{1}}
A spinner is divided into $8$ equal parts: $3$ red, $4$ blue and $1$ yellow.
\begin{center}
\begin{tikzpicture}[font=\footnotesize]
\foreach \c [count=\i from 0] in {tmPink!45,tmBlue!50,tmPink!45,tmBlue!50,StarGold,tmBlue!50,tmPink!45,tmBlue!50}{
  \fill[\c] (0,0) -- (90-45*\i:1.3) arc[start angle=90-45*\i,end angle=45-45*\i,radius=1.3] -- cycle;
  \draw[white,line width=1.2pt] (0,0) -- (90-45*\i:1.3);}
\draw[Slate,line width=0.8pt] (0,0) circle (1.3);
\draw[-{Stealth},line width=1.4pt,Ink] (0,0) -- (25:1.0);
\fill[Ink] (0,0) circle (2.5pt);
\foreach \c/\n [count=\i] in {tmPink!45/Red,tmBlue!50/Blue,StarGold/Yellow}{
  \fill[\c] (2.0,1.0-0.55*\i) rectangle ++(0.32,0.32);
  \node[anchor=west,font=\small] at (2.4,1.16-0.55*\i) {\n};}
\end{tikzpicture}
\end{center}
\pa{a} What is the chance that the arrow stops on blue? \quad \pa{b} About how many times would you expect blue in $80$ spins?
\end{probbox}
\probsol{
\pa{a} $4$ of the $8$ equal parts are blue: $\dfrac48=\dfrac12$.\par
\pa{b} About half of the spins: $80\times\dfrac12=40$ times.
\dapso{\pa{a} $\dfrac12$. \pa{b} About $40$ times.}}
\gap
\begin{probbox}{47}{\stars{1}}
The mean age of a family of four is $24$ years now. What was their mean age $4$ years ago (all four were already born then)?
\end{probbox}
\probsol{
Total age now: $24\times4=96$ years.\par $4$ years ago each person was $4$ years younger: $96-4\times4=80$ years.\par Mean: $80:4=20$ years.
\dapso{$20$ years.}}
\gap
\begin{probbox}{48}{\stars{2}}
In a test, the $32$ pupils of a class have a mean mark of $8$. The $20$ girls have a mean mark of $8.3$. Find the mean mark of the boys.
\end{probbox}
\probsol{
Whole class: $8\times32=256$.\par Girls: $8.3\times20=166$.\par Boys: $32-20=12$ pupils with $256-166=90$ marks.\par Mean: $90:12=7.5$.
\dapso{$7.5$.}}
\gap
\begin{probbox}{49}{\stars{1}}
A box holds $12$ balls, some red and some white. Nam takes a ball, notes its colour and puts it back, $60$ times; red comes up $45$ times. About how many red balls are there in the box?
\end{probbox}
\probsol{
Ratio for red: $\dfrac{45}{60}=\dfrac34$.\par So about $\dfrac34$ of the balls are red: $12\times\dfrac34=9$ balls.
\dapso{About $9$ red balls.}}
\gap
\begin{probbox}{50}{\stars{1}}
An, Bình and Chi stand in a row for a photo.\par
\pa{a} In how many different orders can they stand? \quad \pa{b} In how many of them is An in the middle?
\end{probbox}
\probsol{
\pa{a} First place: $3$ choices, second: $2$, third: $1$; \ $3\times2\times1=6$ orders:\par
ABC, ACB, BAC, BCA, CAB, CBA (A $=$ An, B $=$ Bình, C $=$ Chi).\par
\pa{b} An in the middle: BAC and CAB, i.e.\ $2$ orders.
\dapso{\pa{a} $6$. \pa{b} $2$.}}
